\section{Multi-Step Planning}

\label{sec:planning}

\subsection{Hierarchical Task Decomposition}

Given a high-level goal (e.g., ``Book a flight from SFO to JFK for March 15''), the planner decomposes it into a hierarchy of sub-tasks:

\begin{enumerate}[leftmargin=*]
    \item Open browser and navigate to flight booking site.
    \item Enter departure city (SFO).
    \item Enter destination city (JFK).
    \item Select date (March 15).
    \item Search for flights.
    \item Select preferred flight from results.
    \item Complete booking process.
\end{enumerate}

Each sub-task is further decomposed into atomic actions:

\begin{lstlisting}[language=Python, caption=Task decomposition structure]
@dataclass
class Task:
    goal: str
    subtasks: list["Task"]
    actions: list[Action]  # Leaf level
    status: TaskStatus
    checkpoint: ScreenState | None

    def is_complete(self, state) -> bool:
        """Check completion condition."""

    def rollback_to(self, checkpoint):
        """Restore to checkpoint state."""
\end{lstlisting}

\subsection{Observation-Action Loop}

The core execution loop follows an Observe-Think-Act cycle:

\begin{algorithm}[t]
\caption{Observation-Action Loop}
\label{alg:loop}
\begin{algorithmic}[1]
\REQUIRE Goal $g$, safety policy $\pi$, model $M$
\STATE $\text{plan} \gets M.\text{decompose}(g)$
\STATE $\text{history} \gets []$
\WHILE{not plan.is\_complete()}
    \STATE $s \gets \text{perceive}()$ \COMMENT{Screen understanding}
    \STATE $\text{history.append}(s)$
    \STATE $a \gets M.\text{select\_action}(s, \text{plan}, \text{history})$
    \IF{not $\pi.\text{allows}(a)$}
        \STATE $a \gets M.\text{replan}(s, a, \pi.\text{reason})$
    \ENDIF
    \IF{$a.\text{requires\_confirmation}$}
        \STATE $\text{await\_user\_approval}(a)$
    \ENDIF
    \STATE $\text{result} \gets \text{execute}(a)$
    \STATE $\text{history.append}(a, \text{result})$
    \IF{$\text{result.failed}$}
        \STATE $\text{plan.update}(M.\text{handle\_error}(s, a, \text{result}))$
    \ENDIF
    \STATE $\text{plan.update\_progress}(s)$
\ENDWHILE
\end{algorithmic}
\end{algorithm}

\subsection{Error Recovery}

When actions fail, the planner employs a graduated recovery strategy:

\begin{enumerate}[leftmargin=*]
    \item \textbf{Retry}: For transient failures (element not yet loaded), wait and retry.
    \item \textbf{Alternative action}: Try a different approach to the same sub-goal (e.g., keyboard shortcut instead of menu click).
    \item \textbf{Rollback and replan}: Restore to the last checkpoint and generate a new plan for the remaining sub-tasks.
    \item \textbf{Escalate}: Report failure to the user with context and request guidance.
\end{enumerate}

Checkpoints are created automatically at sub-task boundaries and can be created manually at any point. For browser tasks, checkpoints capture the page URL and relevant form state. For desktop tasks, checkpoints capture a screenshot and the accessibility tree.

\subsection{Context Window Management}

Computer use tasks generate large observation histories that can exceed LLM context windows. Operative manages context through:

\begin{itemize}[leftmargin=*]
    \item \textbf{Summarization}: Older observations are summarized by a secondary LLM call that extracts key state changes and decisions.
    \item \textbf{Relevance filtering}: Only elements relevant to the current sub-task are included in the prompt.
    \item \textbf{Screenshot compression}: Screenshots are resized and compressed before inclusion as vision inputs. Element annotations are overlaid on screenshots to reduce reliance on text descriptions.
    \item \textbf{Sliding window}: The most recent 5 observations are included in full; older observations are represented by their summaries.
\end{itemize}

